% Nouveau EGA V (cote 25): archive cote 25, site page 90 (PDF page 91).
% typescript. Draft compared: none.
% AI-transcribed by Claude Opus 5.5 (claude-opus-5-5, adaptive thinking, effort high), Message Batches API, batch ega25b (msgbatch_01A7n5YuXubUPNK1zsepi8RG), 2026-11-13; checked against both supplied scan parts; see CHECK_DECISIONS.tsv.
\scanpage{25/90}
\mspage{\unsure{84|87}}

\uline{Proposition} 14.5. Supposons dim X $\hand{\leqslant}$ m-1, \var{\xout{Soit Y l'image fermée}}{}{typewriter strike-out}{ega25_25_090-v0001}
\var{\xout{\unsure{et}}}{}{typewriter strike-out}{ega25_25_090-v0002} f:X $\hand{\to}$ P une immersion, \var{\xout{Soit}}{}{typewriter strike-out}{ega25_25_090-v0003} enfin X séparable sur k. Soit Y\hand{$_{\unsure{\eta|y}}$} l'image
schématique de $X_{k(\hand{\unsure{\eta|y}})}$ dans $Q(C_{\hand{\unsure{\eta|y}}})$. Alors le morphisme induit
$p'_{C_{\hand{\unsure{\eta|y}}}} : X_{k(\hand{\unsure{\eta|y}})} \hand{\to} Y_{\hand{\unsure{\eta|y}}}$ est \uline{birationnel}\hand{, et pour tout point x de $X_{k(\unsure{\eta|y})}$ au dessus}
\hand{d'un pt fermé de X, $p'_{C_{\unsure{\eta|y}}}$ est \uline{étale} en x et \unsure{$\overline{m}$|nr} en les points de $p'^{-1}_{C_{\unsure{\eta|y}}}(p'_{C_{\unsure{\eta|y}}}(x))$.}
Signalons la conséquence suivante:

\uline{Corollaire} 14.6. Soit X un schéma algébrique projectif \var{\inserted[typed]{irréductible}}{irréductible}{insertion}{ega25_25_090-v0004} \var{\xout{irréductible}}{}{typewriter strike-out}{ega25_25_090-v0005}
\var{\inserted[typed]{et séparable}}{et séparable}{insertion}{ega25_25_090-v0006}
de dimension n sur un corps infini k. Alors \var{\deleted{X est birationellement}}{}{deletion}{ega25_25_090-v0007} \var{\inserted{$\exists$ un plongement birationnel de X}}{$\exists$ un plongement birationnel de X}{insertion}{ega25_25_090-v0008}
\var{\deleted{isomorphe à}}{}{deletion}{ega25_25_090-v0009} \var{\inserted{\var{\deleted{\unsure{dans}}}{}{deletion}{ega25_25_090-v0011} sur}}{\var{\deleted{\unsure{dans}}}{}{deletion}{ega25_25_090-v0011} sur}{insertion}{ega25_25_090-v0010} une hypersurface \var{\deleted{dans}}{}{deletion}{ega25_25_090-v0012} P$^{n+1}$ .

On se gardera de croire \var{\deleted{que}}{}{deletion}{ega25_25_090-v0013}, même si X est un sous-préschéma fermé
lisse géom irréd de P, de dimension m-1 = n , \var{\inserted[typed]{que}}{que}{insertion}{ega25_25_090-v0014} la projection conique p$_{C_{\hand{\unsure{\eta|y}}}}$
soit \var{\inserted[typed]{nécessairement}}{nécessairement}{insertion}{ega25_25_090-v0015} une \uline{immersion} . En effet, lorsque k est infini, cela impliquerait
qu'il existe une C rationnelle sur k ayant la même propriété, donc que X
est isomorphe à une hypersurface non singulière de \var{\xout{\unsure{P Q}}}{}{typewriter strike-out}{ega25_25_090-v0016} P$^{n+1}$. Or
même pour n=1, (\var{\inserted{\unsure{donc|dans} X}}{\unsure{donc|dans} X }{insertion}{ega25_25_090-v0017}courbe algébrique \var{\inserted[typed]{projective}}{projective}{insertion}{ega25_25_090-v0018} lisse \var{\xout{sur}}{}{typewriter strike-out}{ega25_25_090-v0019} et connexe sur un corps alg.
clos), il est facile de construire des exemples où X ne peut s'immerger
dans un P$^2$ . De même, dans 14.4. , on se gardera de confondre l'énon-%
cé donné avec l'assertion (en général fausse) que p$_C$ soit lui-même un
monomorphisme (contre-exemple précédent\var{\xout{\unsure{), ou que}}}{}{typewriter strike-out}{ega25_25_090-v0020}, ou contre-exemple enco-%
re plus évident si X est lisse de dimension m), ou que p$_C$ soit non
ramifié. Pour ce dernier point, on prendra pour s'en convaincre X
sous-\var{\xout{pré}}{}{typewriter strike-out}{ega25_25_090-v0021}schéma fermé lisse \var{\xout{de dim}}{}{typewriter strike-out}{ega25_25_090-v0022} irréd de dim m (sur k alg clos disons
si \var{\xout{\unsure{X Q}}}{}{typewriter strike-out}{ega25_25_090-v0023} on avait un X $\hand{\to}$ Q $\hand{\simeq}$ P$^m$ non ramifié, il serait étale pour
raisons de dimensions, or on peut montrer (cf Chap VIII !) que cela impli-%
que X $\hand{\xrightarrow{\sim}}$ P$^m$ (P$^m$ étant "simplement connexe). \hand{--} La signification géométri-%
que intuitive de 14.4. est que\var{\xout{\unsure{, pour variable dans}}}{}{typewriter strike-out}{ega25_25_090-v0024} l'ensemble de
ramification de p$_{C_{\hand{\unsure{\eta|y}}}}$ est ``variable'' \var{\inserted{sur k}}{sur k}{insertion}{ega25_25_090-v0025}, plus précisément l'ensemble de
ramification de p$_{C_{\hand{\xi}}}$ pour \hand{\xi} variable dans \var{\inserted{un ouvert de}}{un ouvert de}{insertion}{ega25_25_090-v0026} Grass$^{m+1}$(\hand{$\bar{k}$}) varie dans X(\hand{$\bar{k}$}), et
n'\uline{admet aucun} \var{\deleted{point fermé de X comme}}{}{deletion}{ega25_25_090-v0027} \uline{``point fixe''} ....

\mspage{90}
